\documentclass[11pt]{article}
\usepackage[margin=1in]{geometry}
\usepackage{booktabs}
\title{DME: Lesioned vs. Intact, within SNc --- Interpretation}
\author{GSE233866, 6-OHDA lesion cohort}
\date{}

\begin{document}
\maketitle

\subsection*{Why a pooled network, not the two standalone ones}

Unlike the two standalone SNc networks (\texttt{../snc\_intact\_network/},
\texttt{../snc\_lesioned\_network/}), this analysis pools both conditions'
cells (10{,}928 total: 8{,}637 intact + 2{,}291 lesioned) into a single
network before testing, so module identity is shared and \texttt{FindDMEs}
can compare eigengenes on equal footing rather than comparing across two
independently clustered module sets whose colors and boundaries aren't
guaranteed to line up. This is the general pattern behind every formal DME
test in this project: build one network on the union of the two groups
being compared, then test each module's eigengene between the groups
within that single, shared module definition.

\subsection*{Result}

CACNA1D's module in this pooled network is \textbf{blue} (2{,}249 genes ---
larger than either standalone network's CACNA1D module, since it now
contains cells and genes from both conditions), with $\mathrm{kME}=0.40$.
Its eigengene is \textbf{significantly downregulated after lesioning}
($\mathrm{avg\_log2FC}=-1.66$, $p_{\mathrm{adj}} = 6.1\times10^{-269}$) ---
one of the most statistically extreme results anywhere in this project
(only yellow's upregulation, $p_{\mathrm{adj}}\approx0$, and red's
downregulation, $p_{\mathrm{adj}}=1.1\times10^{-141}$, are in the same
range), and far more extreme than the equivalent human PD-vs-control
comparison, where CACNA1D's own module does not reach significance at all
(\texttt{archive/human\_GSE243639/analysis\_results/pd\_vs\_control\_combined\_dme/}
--- archived 2026-08-13 pending a rebuild from raw data; see
\texttt{archive/human\_GSE243639/README.md} at the repo root).

\subsection*{Reconciling eigengene-level and kME-level results}

Combined with the standalone networks' finding that CACNA1D's own kME
\emph{increases} after lesioning (0.346 $\to$ 0.398), the full picture is:
CACNA1D's broader co-expression program is suppressed in the surviving
lesioned population at the eigengene level (fewer/less-active cells
expressing the program overall), while CACNA1D itself becomes relatively
more central to whatever remains of that shrinking program. These are not
contradictory statistics --- eigengene level and kME are different
quantities (module-wide average activity vs.\ one gene's correlation to
that average) --- but the combination is a specific, testable hypothesis
about what's happening mechanistically, not an established mechanism.
Whether this reflects selective survival of cells with higher CACNA1D
centrality, a genuine transcriptional response to lesioning, or both,
cannot be distinguished from co-expression statistics alone; it would need
an independent perturbation or lineage-tracing experiment to resolve.

\subsection*{Caveats}

The brown module's log2FC is flagged unstable in the source table (module
eigengenes can be negative or near-zero, which destabilizes a ratio-based
fold-change); its $p_{\mathrm{adj}}=7.1\times10^{-17}$ is unaffected by that
instability and should still be trusted for direction and significance.
More generally: with 2{,}249 genes, blue is the single largest module in
this project, so its enrichment profile (not reported in this file; see
\texttt{enrichr\_table.csv} directly) is likely to be broad and less
specific than the smaller, more tightly-themed modules in the standalone
networks --- a large pooled module absorbs genes from what might be several
more specific sub-programs in a smaller, single-condition network.

\end{document}
